%! Introducing Statistics: Why Be Normal? — Unit 6, Lesson 6.1 — Homework (Answer Key)
\documentclass[10pt]{article}
\usepackage{apstats-article}
\usepackage{apstats-key}   % pulls in -boxes; do NOT also load -boxes
\usepackage{grading}
\gradingassignment{apstats-u06-l01-hw}
\gradingcourse{advanced_statistics}
\newcommand{\TallMath}[1]{$\displaystyle #1\rule[-1.4em]{0pt}{3.2em}$}

\begin{document}

\pageheader{Unit 6, Lesson 6.1}{Homework --- Answer Key}

\vspace{0.15in}

\begin{enumerate}

  \item For each scenario, (a) verify the Large Counts condition, (b) calculate $SE_{\hat p}$,
        and (c) identify the inference goal as \emph{confidence interval} (CI) or
        \emph{significance test} (Test).

        \begin{enumerate}[label=\alph*.]
          \item[\textbf{I.}] A poll of 350 randomly selected registered voters finds $\hat p = 0.58$
                support a ballot measure. Researchers want to estimate the true proportion
                of all registered voters who support it.
                \begin{itemize}
                  \item Large Counts: \hfill Met? \smash{\gradedblank[id=hw1a2, type=short, points=1, answer=yes, check=text]{1.5cm}{0.17in}}\gradedhide{\\[4pt]\gradedhere{\gradedblank[id=hw1a1, type=rubric, points=1, display={\bfseries $203 \ge 10$ and $147 \ge 10$}]{\linewidth+6pt}{\gradedlinesheight{9}}}\writelines{10}}
                  \item $SE_{\hat p} = $ \smash{\gradedblank[id=hw1a3, type=short, points=1, answer=0.0264, check=numeric, tolerance=0.0005]{6cm}{0.17in}}
                  \item Inference goal: \gradedblank[id=hw1a4, type=choice, points=1, answer=CI, display={CI}]{4cm}{0.17in}
                \end{itemize}

          \newpage
          \item[\textbf{II.}] A nutritionist claims that 30\% of teenagers eat breakfast daily.
                A researcher surveys 90 randomly selected teenagers and finds $\hat p = 0.24$.
                She wants to test whether the true proportion is less than 30\%.
                \begin{itemize}
                  \item Large Counts (use $p_0 = 0.30$): \hfill Met? \smash{\gradedblank[id=hw1b2, type=short, points=1, answer=yes, check=text]{1.5cm}{0.17in}}\gradedhide{\\[4pt]\gradedhere{\gradedblank[id=hw1b1, type=rubric, points=1, display={\bfseries $27 \ge 10$ and $63 \ge 10$}]{\linewidth+6pt}{\gradedlinesheight{9}}}\writelines{10}}
                  \item $SE_{\hat p} = \sqrt{\hat p(1-\hat p)/n} = $ \smash{\gradedblank[id=hw1b3, type=short, points=1, answer=0.0450, check=numeric, tolerance=0.0005]{5cm}{0.17in}}
                  \item Inference goal: \gradedblank[id=hw1b4, type=choice, points=1, answer=Test, display={Significance Test}]{4cm}{0.17in}
                \end{itemize}

          \item[\textbf{III.}] A manufacturing plant samples 500 units and finds $\hat p = 0.03$
                are defective. Quality control wants to report a range for the true defect rate.
                \begin{itemize}
                  \item Large Counts: \hfill Met? \smash{\gradedblank[id=hw1c2, type=short, points=1, answer=yes, check=text]{1.5cm}{0.17in}}\gradedhide{\\[4pt]\gradedhere{\gradedblank[id=hw1c1, type=rubric, points=1, display={\bfseries $15 \ge 10$ and $485 \ge 10$}]{\linewidth+6pt}{\gradedlinesheight{9}}}\writelines{10}}
                  \item $SE_{\hat p} = $ \smash{\gradedblank[id=hw1c3, type=short, points=1, answer=0.0076, check=numeric, tolerance=0.0005]{6cm}{0.17in}}
                  \item Inference goal: \gradedblank[id=hw1c4, type=choice, points=1, answer=CI, display={CI}]{4cm}{0.17in}
                \end{itemize}
        \end{enumerate}

  \boxguard[12]
  \item Complete the table comparing confidence intervals and significance tests.

        \renewcommand{\arraystretch}{2.0}
        \begin{tabular}{@{} >{\bfseries}p{0.24\linewidth} p{0.34\linewidth} p{0.34\linewidth} @{}}
          & \textbf{Confidence Interval} & \textbf{Significance Test} \\
          \hline
          Research goal & \ansline{Estimate the unknown parameter}\gradedkeyline{\gradedblank[id=hw2a, type=short, points=1, answer={estimate the parameter}, check=text, display={Estimate the parameter}]{\linewidth+6pt}{0.3in}} & \ansline{Test a specific claimed parameter value}\gradedkeyline{\gradedblank[id=hw2b, type=short, points=1, answer={test a claimed parameter value}, check=text, display={Test a claimed parameter value}]{\linewidth+6pt}{0.3in}} \\
          Key output    & \ansline{A range of plausible parameter values}\gradedkeyline{\gradedblank[id=hw2c, type=short, points=1, answer={a range of plausible values}, check=text, display={A range of plausible values}]{\linewidth+6pt}{0.3in}} & \ansline{A $p$-value and reject/fail-to-reject decision}\gradedkeyline{\gradedblank[id=hw2d, type=short, points=1, answer={a p-value and a decision}, check=text, display={A $p$-value and a decision}]{\linewidth+6pt}{0.3in}} \\
          Based on      & \multicolumn{2}{l}{Normal sampling distribution of $\hat p$ (when conditions met)} \\
        \end{tabular}

\end{enumerate}

\begin{reflectionbox}
\textbf{Preview:} Lesson 6.2 introduces the one-sample $z$-interval for a population
proportion --- the first complete inference procedure. Review the formula
$\hat p \pm z^* \sqrt{\hat p(1-\hat p)/n}$ and think about what each part represents.
\end{reflectionbox}

\end{document}
